\chapter{方法：Texture Resonance Retrieval}

TRR 的核心思想是用二阶统计描述音频纹理。给定深层音频特征矩阵 $\mathbf{H}\in\mathbb{R}^{C\times T}$，其中 $C$ 表示通道维度，$T$ 表示帧数，TRR 计算归一化 Gram 矩阵
\[
	\mathbf{G}=\frac{1}{T}\mathbf{H}\mathbf{H}^{\top}, \qquad
	\mathbf{g}=\frac{\mathrm{vec}(\mathbf{G})}{\|\mathrm{vec}(\mathbf{G})\|_2}.
\]
检索时，查询音频和数据库音频都被映射为 $\mathbf{g}$，然后用余弦相似度排序。该表示保留跨通道共激活结构，弱化绝对时间对齐；这与吉他效果中的失真颗粒、调制纹理、空间尾音和动态密度等现象相符，因为这些现象常表现为稳定纹理，而非某个精确时间点的瞬时事件。

与 mean pooling 相比，Gram 统计的优势在于它保留了通道之间的二阶关系。若两个音频片段在整体均值上相近，但一个包含更强的调制共振或失真相关结构，一阶摘要可能无法区分，而二阶统计可能改变其邻域排序。当前实验结果不能证明这个假设在所有音频任务上成立，但它支持该假设在 guitar-effects 参数检索任务上的有效性。

\begin{figure}[H]
	\centering
	\resizebox{0.92\textwidth}{!}{%
	\begin{tikzpicture}[node distance=1.35cm and 1.35cm]
		\node[reportnode, minimum width=3.0cm] (audio) {参考音频片段};
		\node[reportprocess, right=of audio, minimum width=3.2cm] (feat) {深层帧特征\\$\mathbf{H}\in\mathbb{R}^{C\times T}$};
		\node[reportprocess, right=of feat, minimum width=3.0cm] (gram) {Gram 统计\\$\mathbf{H}\mathbf{H}^{\top}/T$};
		\node[reportprocess, right=of gram, minimum width=3.0cm] (vec) {向量化与归一化\\$\mathbf{g}$};
		\node[reportnode, below=1.0cm of vec, minimum width=3.2cm] (knn) {余弦近邻检索\\Top-1/Top-K};
		\node[reportnode, left=of knn, minimum width=3.1cm] (preset) {可执行预设参数\\$\hat{\theta}$};
		\draw[reportarrow] (audio) -- (feat);
		\draw[reportarrow] (feat) -- (gram);
		\draw[reportarrow] (gram) -- (vec);
		\draw[reportarrow] (vec) -- (knn);
		\draw[reportarrow] (knn) -- (preset);
	\end{tikzpicture}}
	\caption{TRR 的计算路径。当前实验主要使用缓存 TRR 向量；缺失缓存时才懒加载编码器计算嵌入。}
	\label{fig:trr-mechanism}
\end{figure}

图~\ref{fig:trr-mechanism} 也说明了 TRR 的边界。TRR 只是检索表示，不负责生成新音频，也不直接修复插件参数。当前主实验的统计证据链只覆盖检索式参数对齐；未进入该证据链的系统探索不应被写成已验证贡献。

P0 E3 之后，fusion 不再只是辅助诊断机制，而是论文第二条实验证据线。它根据文本侧和音频侧 top-k 分数计算熵、质量分数和 audio bias，再输出 $w_{\text{text}}$ 与 $w_{\text{audio}}$。该机制的价值在于真实退化下避免 audio-heavy fixed fusion 的脆弱性，并把系统推向更稳定的中高文本权重工作区间。它的局限也很明确：当文本与音频发生语义冲突时，旧 Protocol-C 仍显示当前规则可能偏离更低误差的 TRR-only 结果。因此，本文把 adaptive fusion 写成当前退化集合下的鲁棒机制，不写成跨域或冲突场景下的通用最优策略。
